\begin{textsample}{POS Dim 2 – ms – Score 19.00 – ms/ms\_em\_38.txt}  \label{ex:f2_pos_050}
4.3.4.
Unidades Curriculares \textbf{Eletivas} A Unidade Curricular \textbf{Eletiva} compõe a estrutura do Aprofundamento em Área de Conhecimento dos \textbf{Itinerários} \textbf{Formativos} e permite o conhecimento de diferentes temas, vivências e aprendizagens.
Esse componente ocupa um lugar de destaque na diversificação das experiências escolares e oferece um espaço privilegiado para a experimentação, a interdisciplinaridade e o aprofundamento dos estudos.
Assim, é possível propiciar o desenvolvimento de diferentes linguagens – verbal, musical, matemática, gráfica, plástica, corporal, visual –, além de consolidar competências previstas neste \textbf{Currículo}.
A elaboração da Unidade Curricular \textbf{Eletiva} é responsabilidade da escola, e deve possibilitar a experimentação em diferentes temas, vivências e aprendizagens, ancorada em um trabalho pedagógico intencional, estruturado, com a participação ativa dos estudantes, e pautada na flexibilização, na criatividade e na interdisciplinaridade.
Com forte presença de atividades lúdicas, a intencionalidade pedagógica deve ser clara e articulada com as Áreas de Conhecimento, os Eixos \textbf{Estruturantes} e as Competências fundamentais.
É fundamental que o professor participe do processo de repensar a construção do conhecimento, na qual a mediação e a interação são os pressupostos essenciais para que ocorra a aprendizagem.
Para esse fim, o uso de metodologias ativas, como forma de desenvolver o processo do aprender, envolve estudantes e professores de forma conjunta na busca da (re)construção dos saberes, de forma a propiciar, além da aprendizagem cognitiva, a formação de valores para a vida.
A opção por uma metodologia ativa deve considerar que os estudantes desenvolvam mecanismos de problematização, para que tenham a possibilidade de examinar, refletir, posicionar -se de forma crítica, aprender a expor sua opinião e a respeitar pensamentos diferentes ; dessa forma, o planejamento deve ser pautado nas necessidades e interesses dos estudantes, por meio de práticas pedagógicas que os envolvam no próprio processo de aprendizagem. 4.4.
Estrutura de Funcionamento da Flexibilização Curricular A organização da parte \textbf{flexível} do \textbf{currículo} embasou -se na Lei Federal n. 13.415/2017, na Resolução CNE/CEB n. 03/2018, na \textbf{Portaria} MEC n. 1.432/2018 e no \textbf{Guia} de Implementação do Novo Ensino Médio, considerando as especificidades de cada município/região do Estado.
Assim, a composição do Aprofundamento em Área de Conhecimento dos \textbf{Itinerários} \textbf{Formativos} torna -se responsabilidade das escolas, haja vista que o produto final deve \textbf{atender} às aspirações e características do contexto local.
Nesta proposta, a SED/MS elaborou as UCs, articuladas aos arranjos e potenciais locais, às demandas e necessidades do mundo contemporâneo e ao interesse dos estudantes, organizando -as por área de conhecimento no \textbf{Catálogo} de Unidades Curriculares.
O \textbf{Catálogo} reúne informações detalhadas acerca de cada UC, tais como : nome, objetivo/descrição, \textbf{carga} \textbf{horária}, competências e habilidades a serem desenvolvidas, relação de aprofundamento com a BNCC, objetos de conhecimento a serem mobilizados, sugestões didáticas, eixo \textbf{estruturante} principal, eixos \textbf{estruturantes} secundários, estrutura necessária para a \textbf{oferta}, \textbf{perfil} do professor, dentre outras.
Analisado o \textbf{Catálogo}, as escolas procederão à composição dos \textbf{Itinerários} \textbf{Formativos}, os quais serão \textbf{ofertados} em turmas organizadas por grupos de interesse e não por agrupamentos \textbf{seriados}.
Cada escola, mediante escuta à comunidade escolar, a fim de coletar informações a respeito das principais áreas de interesse, fará a composição de, no mínimo, dois \textbf{Itinerários} \textbf{Formativos}, em áreas distintas, de modo a garantir o estabelecido no parágrafo 6º do artigo 12 da Resolução CNE/CEB n. 03/2018.
Ressalta -se que os esforços envidados para a composição dos \textbf{Itinerários} \textbf{Formativos} devem garantir a efetiva possibilidade de escolha e protagonismo dos estudantes, atentando -se para o fato de que não podem ser repetições ou mero reforço das competências e habilidades mobilizadas na Formação Geral Básica, mas uma ampliação em suas respectivas Áreas, na Formação \textbf{Profissional} e \textbf{Técnica} ou em Áreas Integradas, de tal forma que os estudantes tenham a oportunidade de escolher o Aprofundamento desejado.

% matched lemmas: atender, carga, catálogo, currículo, eletivo, estruturante, flexível, formativo, guia, horário, itinerário, oferta, ofertar, perfil, portaria, profissional, seriar, técnico
\end{textsample}
